% 数学ノート集・共通テンプレート v2（積分論2026年度第1部の版面を参考）
% A5判。章内で定義・定理・例・注意等の番号を共通化。
% コンパイル：platex main.tex（3回），dvipdfmx main.dvi
%              dvipdfmx math_note_template_v2.dvi
\documentclass[a5paper,10pt,onecolumn]{jsarticle}
\usepackage[dvipdfmx,top=22mm,bottom=20mm,inner=19mm,outer=18mm,headheight=13pt,headsep=6mm]{geometry}
\usepackage{amsmath,amssymb,amsthm,mathtools}
\usepackage{enumitem,fancyhdr,titlesec}
\usepackage[dvipdfmx,bookmarks=true,colorlinks=false]{hyperref}
\usepackage{pxjahyper}

% ----- 差し替える文書情報 -----
\newcommand{\NoteTitle}{Hahn--Banachの定理と凸集合の分離}
\newcommand{\NoteSeries}{数学ノート集}

% ----- 文字・組版 -----
% 参考PDFの質素な明朝＋Computer Modernの印象を引き継ぎながら，
% A5判で窮屈にならないよう本文の行間を少し確保する．
\linespread{1.13}
\setlength{\parindent}{1zw}
\setlength{\parskip}{0pt}
\setlength{\abovedisplayskip}{6pt plus 2pt minus 2pt}
\setlength{\belowdisplayskip}{6pt plus 2pt minus 2pt}
\setlength{\abovedisplayshortskip}{3pt plus 2pt minus 1pt}
\setlength{\belowdisplayshortskip}{4pt plus 2pt minus 1pt}
\setlist[enumerate]{topsep=3pt,itemsep=1pt,parsep=0pt,leftmargin=2.4zw}
\titleformat{\section}{\normalfont\normalsize\bfseries}{\thesection}{0.8zw}{}
\titlespacing*{\section}{0pt}{2.0ex plus .4ex}{.9ex}
\titleformat{\subsection}{\normalfont\normalsize\bfseries}{\thesubsection}{0.7zw}{}
\titlespacing*{\subsection}{0pt}{1.5ex plus .3ex}{.6ex}
\raggedbottom

% ----- 柱とページ番号：ページ上部、罫線なし -----
\fancypagestyle{mathnote}{%
  \fancyhf{}%
  \fancyhead[L]{\footnotesize \NoteSeries\ （\NoteTitle）}%
  \fancyhead[R]{\footnotesize \thepage}%
  \renewcommand{\headrulewidth}{0pt}%
  \renewcommand{\footrulewidth}{0pt}%
}
\pagestyle{mathnote}

% ----- 節内通し番号の定義・定理群 -----
\newtheoremstyle{mathnotetheorem}
 {4.5pt}{4.5pt}{\normalfont}{}%
 {\bfseries}{}{0.65zw}%
 {\thmname{#1}\thmnumber{ #2}\thmnote{（#3）}}
\theoremstyle{mathnotetheorem}
\newtheorem{theorem}{定理}[section]
\newtheorem{definition}[theorem]{定義}
\newtheorem{lemma}[theorem]{補題}
\newtheorem{proposition}[theorem]{命題}
\newtheorem{corollary}[theorem]{系}
\newtheorem{example}[theorem]{例}
\newtheorem{remark}[theorem]{注意}
\renewcommand{\proofname}{証明}
\renewcommand{\qedsymbol}{$\square$}
\makeatletter
\renewenvironment{proof}[1][\proofname]{%
  \par\pushQED{\qed}\normalfont\topsep4\p@\@plus2\p@\relax
  \trivlist\item[\hskip\labelsep\bfseries #1]\ignorespaces
}{\popQED\endtrivlist\@endpefalse}
\makeatother

% ----- 使うノートで適宜拡張する記号 -----
\newcommand{\R}{\mathbb{R}}
\newcommand{\norm}[1]{\lVert #1\rVert}
\newcommand{\abs}[1]{\lvert #1\rvert}

\begin{document}
\section*{読了後にできること}
本稿を読むための前提は，線形汎関数，作用素ノルム，Hahn--Banachの定理．
本稿の具体的な到達課題は，部分空間の汎関数をノルムを保って延長する．

\section{線形汎関数の拡張問題}

本稿では $\mathbb K=\mathbb R$ または $\mathbb C$ とする．
線形空間の部分空間 $M\subset X$ 上に定義された線形汎関数を，
定められた評価を保ちつつ $X$ 全体へ拡張する問題を考える．
実数体の場合には劣線形汎関数が基本的な役割を果たす．

\begin{definition}
実線形空間 $X$ 上の関数 $p:X\to\mathbb R$ が
\emph{劣線形}であるとは
\begin{align*}
 p(x+y)&\le p(x)+p(y)\quad(x,y\in X),\\
 p(tx)&=t\,p(x)\quad(x\in X,\ t\ge0)
\end{align*}
を満たすことをいう．
正値性 $p(x)\ge0$ は要求しない．
\end{definition}

\begin{lemma}\label{lem:sublinear}
劣線形汎関数 $p$ は $p(0)=0$ と
$-p(-x)\le p(x)$ を満たす．
また実線形汎関数 $f:X\to\mathbb R$ が
$f\le p$ を満たすなら
\[
 -p(-x)\le f(x)\le p(x)\qquad(x\in X)
\]
である．
\end{lemma}
\begin{proof}
正斉次性に $t=0$ を代入すれば $p(0)=0$ である．
劣加法性により
$0=p(0)=p(x+(-x))\le p(x)+p(-x)$ となる．
$f\le p$ なら $f(-x)\le p(-x)$ であり，
$f(-x)=-f(x)$ を用いて左側の不等式を得る．
\end{proof}

\begin{example}
ノルム空間のノルム $p(x)=\|x\|$ は実線形空間上の
劣線形汎関数である．また任意の実線形汎関数
$\varphi:X\to\mathbb R$ それ自体も劣線形である．
\end{example}
\begin{proof}
ノルムの場合は三角不等式と正の斉次性を用いればよい．
線形汎関数の場合は
$\varphi(x+y)=\varphi(x)+\varphi(y)$，
$\varphi(tx)=t\varphi(x)$ が成り立つから
劣線形の両条件を満たす．
\end{proof}

\section{実Hahn--Banachの定理}

\begin{lemma}[1次元拡張補題]\label{lem:onedim}
$X$ を実線形空間，$p:X\to\mathbb R$ を劣線形汎関数，
$M\subset X$ を線形部分空間とする．
$f:M\to\mathbb R$ が線形で $f(m)\le p(m)$
をすべての $m\in M$ について満たすとする．
$v\in X\setminus M$ を取るとき，
$f$ を $M+\mathbb Rv$ 上の線形汎関数 $F$ に拡張して
$F\le p$ を保つことができる．
\end{lemma}
\begin{proof}
$M+\mathbb Rv$ の各元は一意に $m+tv$ と書ける．
$F(m+tv)=f(m)+tc$ と置くことを考える．
$t>0$ のときに $F(m+tv)\le p(m+tv)$ となるためには
\[
 c\le p(m+v)-f(m)\qquad(m\in M)
\]
が必要十分であり，$t<0$ のときの条件は
\[
 f(m)-p(m-v)\le c\qquad(m\in M)
\]
と同値である．従って
\[
 \sup_{m\in M}\{f(m)-p(m-v)\}
 \le c\le
 \inf_{n\in M}\{p(n+v)-f(n)\}
 \tag{$*$}
\]
を満たす実数 $c$ を選べばよい．

任意の $m,n\in M$ について
\[
 f(m)+f(n)=f(m+n)\le p(m+n)
 \le p(m-v)+p(n+v)
\]
だから
\[
 f(m)-p(m-v)\le p(n+v)-f(n).
\]
従って $(*)$ の左端は右端以下である．
しかも $n=0$ を固定すると左端は
$p(v)$ 以下であり，$m=0$ を固定すると右端は
$-p(-v)$ 以上である．両端は有限な実数であるため，
$(*)$ を満たす $c\in\mathbb R$ が存在する．
この $c$ によって定めた $F$ は線形であり，
上記の $t>0,t<0$ の議論から $F\le p$ となる．
$t=0$ ではもとの仮定による．
\end{proof}

\begin{theorem}[実Hahn--Banachの定理]\label{thm:hb-real}
$X$ を実線形空間，$M$ を線形部分空間，
$p:X\to\mathbb R$ を劣線形汎関数とする．
線形汎関数 $f:M\to\mathbb R$ が
$f(m)\le p(m)$（$m\in M$）を満たすなら，
ある線形汎関数 $F:X\to\mathbb R$ が存在して
\[
 F|_M=f,\qquad F(x)\le p(x)\quad(x\in X)
\]
を満たす．
\end{theorem}
\begin{proof}
$M$ を含む部分空間 $N\subset X$ と，
$N$ 上の線形汎関数 $g$ の組で
$g|_M=f$，$g\le p|_N$ を満たすもの全体を
$\mathcal P$ とする．
$(N_1,g_1)\preceq(N_2,g_2)$ を
$N_1\subset N_2$ かつ $g_2|_{N_1}=g_1$
によって定めると，$\mathcal P$ は順序集合となり，
$(M,f)$ を含むため空でない．

$\mathcal C\subset\mathcal P$ を全順序部分集合とする．
$\mathcal C=\varnothing$ なら $(M,f)$ が上界である．
空でない場合には
\[
 N_\mathcal C=\bigcup_{(N,g)\in\mathcal C}N
\]
は線形部分空間である．実際，
$x\in N_1$, $y\in N_2$ のとき，
全順序性により $N_1\subset N_2$ または
$N_2\subset N_1$ であるから，
$x+y$ とスカラー倍は共に $N_\mathcal C$ に属する．
$G(x)=g(x)$（$x\in N$）と定めると，
全順序性と拡張条件によって $G$ は矛盾なく定義され，
線形であり，$G|_M=f$ かつ $G\le p$ を満たす．
従って $(N_\mathcal C,G)$ は $\mathcal C$ の上界である．

Zornの補題から $\mathcal P$ には極大元 $(N_0,g_0)$ が存在する．
$N_0\ne X$ なら $v\in X\setminus N_0$ を取れる．
補題\ref{lem:onedim}により $g_0$ は
$N_0+\mathbb Rv$ に条件を保って拡張され，
極大性に矛盾する．
よって $N_0=X$ であり，$F=g_0$ が求める拡張である．
\end{proof}

\begin{remark}
定理\ref{thm:hb-real}には $X$ の位相もノルムも必要ない．
全空間への拡張の段階でZornの補題を使用する．
後述するノルム空間上の定理では，
劣線形汎関数としてノルムの定数倍を選ぶ．
\end{remark}

\section{ノルムを保存する拡張}

\begin{theorem}[実ノルム空間のHahn--Banachの定理]
$X$ を実ノルム空間，$M\subset X$ を線形部分空間，
$f\in M^*$ とする．
このとき $F\in X^*$ で
\[
 F|_M=f,\qquad \|F\|_{X^*}=\|f\|_{M^*}
\]
を満たすものが存在する．
\end{theorem}
\begin{proof}
$p(x)=\|f\|\|x\|$ と置くと $p$ は劣線形であり，
$f(m)\le|f(m)|\le p(m)$ を満たす．
定理\ref{thm:hb-real}から
$F:X\to\mathbb R$ で $F|_M=f$，
$F(x)\le\|f\|\|x\|$ を満たすものが存在する．
$-x$ にも適用すると
$-F(x)\le\|f\|\|x\|$ であるから
$|F(x)|\le\|f\|\|x\|$ となる．
従って $\|F\|\le\|f\|$ である．
一方 $F$ は $f$ の拡張なので
$\|F\|\ge\|f\|$ が従い，等号となる．
\end{proof}

\section{複素Hahn--Banachの定理}

\begin{lemma}\label{lem:realpart}
$X$ を複素線形空間とし，$f:X\to\mathbb C$
を複素線形汎関数とする．
$u=\operatorname{Re}f$ と置けば，
$u$ は $X$ を実線形空間とみたときの実線形汎関数であり，
\[
 f(x)=u(x)-i\,u(ix)\qquad(x\in X)
\]
が成り立つ．
逆に任意の実線形汎関数 $u:X\to\mathbb R$ に対し，
この式で定義した $f$ は複素線形で
$\operatorname{Re}f=u$ を満たす．
\end{lemma}
\begin{proof}
$f(ix)=if(x)$ なので
$u(ix)=\operatorname{Re}(if(x))=-\operatorname{Im}f(x)$ となり，
最初の式を得る．
逆に $f(x)=u(x)-iu(ix)$ と定義すると
実数についての線形性は明らかであり，
\[
 f(ix)=u(ix)-iu(i^2x)
 =u(ix)+iu(x)=if(x)
\]
となる．従って実線形性と $f(ix)=if(x)$ により
複素線形性が従う．実部は定義から $u$ である．
\end{proof}

\begin{theorem}[複素ノルム空間のHahn--Banachの定理]\label{thm:hb-complex}
$X$ を複素ノルム空間，$M\subset X$ を
複素線形部分空間，$f\in M^*$ とする．
このとき $F\in X^*$ が存在して
\[
 F|_M=f,\qquad \|F\|=\|f\|
\]
が成り立つ．
\end{theorem}
\begin{proof}
$X$ を実ノルム空間とみなし，
$u=\operatorname{Re}f$ を $M$ 上の実線形汎関数とする．
$u(m)\le|f(m)|\le\|f\|\|m\|$ だから
実Hahn--Banachの定理により
実線形汎関数 $U:X\to\mathbb R$ が存在して
$U|_M=u$ かつ
$U(x)\le\|f\|\|x\|$ となる．
\[
 F(x)=U(x)-iU(ix)
\]
と置くと，補題\ref{lem:realpart}から $F$ は複素線形であり，
$M$ 上では $f$ に一致する．

$F(x)\ne0$ とし，
$|\zeta|=1$ かつ $\zeta F(x)=|F(x)|$ を満たす
複素数 $\zeta$ を取る．複素線形性から
\[
 |F(x)|=F(\zeta x)=\operatorname{Re}F(\zeta x)
 =U(\zeta x)\le\|f\|\,\|\zeta x\|
 =\|f\|\,\|x\|.
\]
$F(x)=0$ のときも同じ評価は自明である．
従って $\|F\|\le\|f\|$ となり，
拡張であることから逆向きの不等式も得る．
\end{proof}

\begin{remark}
実数体・複素数体のいずれでも，ノルムを保存する
Hahn--Banachの拡張に完備性は必要ない．
一方，複素数体の証明では実部を用いて実の場合へ帰着するため，
内積空間における複素共役の規約に依存しない．
\end{remark}

\section{双対空間が点を分離すること}

\begin{theorem}[ノルムを与える汎関数]\label{thm:norming}
$X$ を $\mathbb K$ 上のノルム空間，
$x_0\in X\setminus\{0\}$ とする．
このときある $f\in X^*$ が存在して
\[
 \|f\|=1,\qquad f(x_0)=\|x_0\|
\]
を満たす．
\end{theorem}
\begin{proof}
$M=\mathbb Kx_0$ とし，
$f_0(ax_0)=a\|x_0\|$ と定める．
$x_0\ne0$ なのでこれはwell-definedな
$\mathbb K$ 線形汎関数であり，
$|f_0(ax_0)|=|a|\|x_0\|=\|ax_0\|$ により
$\|f_0\|=1$ である．
実数体の場合には実ノルム空間の拡張定理を，
複素数体の場合には定理\ref{thm:hb-complex}を適用し，
$\|f\|=1$ を保つ拡張 $f$ を得る．
\end{proof}

\begin{corollary}[ノルムの双対表示]\label{cor:dualnorm}
任意のノルム空間 $X$ と $x\in X$ に対し
\[
 \|x\|=
 \sup\{|f(x)|:f\in X^*,\ \|f\|\le1\}
\]
が成り立つ．
特に $x\ne y$ なら $f(x)\ne f(y)$ を満たす
$f\in X^*$ が存在する．
\end{corollary}
\begin{proof}
任意の $\|f\|\le1$ に対して
$|f(x)|\le\|x\|$ なので右辺は $\|x\|$ 以下である．
$x\ne0$ なら定理\ref{thm:norming}により
$f(x)=\|x\|$ となる $\|f\|=1$ の汎関数が存在する．
$x=0$ は自明である．
最後に $x-y\ne0$ に定理\ref{thm:norming}を適用する．
\end{proof}

\begin{proposition}[第二双対への埋め込み]\label{prop:bidual}
ノルム空間 $X$ に対し，
\[
 J:X\to X^{**},\qquad (Jx)(f)=f(x)
\]
と定めると，$J$ は線形な等距離埋め込みである．
\end{proposition}
\begin{proof}
各 $x\in X$ について
$|Jx(f)|=|f(x)|\le\|f\|\|x\|$ だから
$Jx\in X^{**}$ で $\|Jx\|\le\|x\|$ である．
$J$ の線形性は $f$ の線形性から従う．
系\ref{cor:dualnorm}によって
$\|Jx\|=\sup_{\|f\|\le1}|f(x)|=\|x\|$ となる．
従って $J$ は等距離であり，特に単射である．
\end{proof}

\begin{remark}
$J:X\to X^{**}$ が全射であるとは限らない．
全射となるとき $X$ は反射的であるという．
Hilbert空間の反射性はRieszの表現定理から導けるが，
一般のBanach空間に対しては成立しない．
\end{remark}

\section{閉部分空間と距離の双対表示}

\begin{lemma}[商ノルム]
$X$ をノルム空間，$M\subset X$ を閉線形部分空間とする．
商空間 $X/M$ に
\[
 \|x+M\|_{X/M}=\inf_{m\in M}\|x-m\|
\]
を定めると，これはノルムになる．
\end{lemma}
\begin{proof}
代表元を $x$ から $x+m_0$ に替えても
下限は変わらない．
斉次性は $M$ の線形性から従い，
三角不等式は
\[
 \|x+y-(m+n)\|\le\|x-m\|+\|y-n\|
\]
で $m,n\in M$ に関する下限を取れば得られる．
またノルムが $0$ なら $x\in\overline M=M$ である．
逆に $x\in M$ なら下限は $0$ である．
\end{proof}

\begin{theorem}[閉部分空間からの分離]\label{thm:subspace-separate}
$M$ をノルム空間 $X$ の閉線形部分空間とし，
$x_0\in X\setminus M$ とする．
このとき $f\in X^*$ で
\[
 \|f\|=1,\qquad f|_M=0,\qquad
 f(x_0)=d(x_0,M)>0
\]
を満たすものが存在する．
\end{theorem}
\begin{proof}
商空間 $X/M$ はノルム空間であり，
$x_0+M\ne0$ である．
定理\ref{thm:norming}を $x_0+M$ に適用し，
$g\in(X/M)^*$ で
\[
 \|g\|=1,\qquad
 g(x_0+M)=\|x_0+M\|=d(x_0,M)
\]
を満たすものを取る．
自然な商写像 $\pi:X\to X/M$ を用いて
$f=g\circ\pi$ と置くと $f$ は線形，
$f|_M=0$ であり，
\[
 |f(x)|\le\|g\|\|\pi x\|
 \le\|x\|
\]
なので $\|f\|\le1$ となる．
一方 $f(x_0)=d(x_0,M)$ だから
任意の $m\in M$ に対して
\[
 d(x_0,M)=|f(x_0-m)|
 \le\|f\|\,\|x_0-m\|.
\]
$m$ に関する下限を取ると
$d(x_0,M)\le\|f\|\,d(x_0,M)$ を得る．
$M$ は閉で $x_0\notin M$ だから距離は正であり，
$\|f\|\ge1$ である．従って $\|f\|=1$ となる．
\end{proof}

\begin{corollary}[距離の双対表示]\label{cor:distance}
$M\subset X$ を閉線形部分空間とする．
任意の $x\in X$ に対して
\[
 d(x,M)=
 \sup\{|f(x)|:f\in X^*,\
           \|f\|\le1,\ f|_M=0\}
\]
が成り立つ．
\end{corollary}
\begin{proof}
$M$ 上で $0$ となる $\|f\|\le1$ の汎関数について
任意の $m\in M$ に対して
$|f(x)|=|f(x-m)|\le\|x-m\|$ である．
下限を取ると右辺は $d(x,M)$ 以下となる．
$x\in M$ なら両辺は $0$ である．
$x\notin M$ なら定理\ref{thm:subspace-separate}の
$f$ が等号を実現する．
\end{proof}

\section{Minkowski汎関数}

\begin{definition}
実ノルム空間 $X$ の部分集合 $C$ が\emph{吸収的}であるとは，
任意の $x\in X$ に対してある $t>0$ が存在して
$x\in tC$ となることをいう．
$C$ が空でない凸集合で $0\in C$ かつ吸収的なら，
\[
 p_C(x)=\inf\{t>0:x\in tC\}
\]
を $C$ の\emph{Minkowski汎関数}という．
\end{definition}

\begin{lemma}\label{lem:gauge}
$C$ を実ノルム空間 $X$ の開凸集合で $0\in C$ とする．
このとき $C$ は吸収的であり，
$p_C:X\to[0,\infty)$ は劣線形である．
さらに
\[
 C=\{x\in X:p_C(x)<1\}
\]
が成り立ち，ある $r>0$ に対して
\[
 p_C(x)\le\frac{\|x\|}{r}\qquad(x\in X)
\]
が成り立つ．
\end{lemma}
\begin{proof}
$C$ は $0$ の開近傍であるから
$B_X(0,r)\subset C$ となる $r>0$ が存在する．
$x\ne0$ なら $t>\|x\|/r$ に対して
$x/t\in B_X(0,r)\subset C$ であり，
$x\in tC$ となる．よって吸収的であり，
$p_C(x)\le\|x\|/r$ が成り立つ．$x=0$ も明らかである．

$s>0$ に対して
\[
 p_C(sx)=\inf\{t>0:sx\in tC\}
 =s\inf\{u>0:x\in uC\}=s\,p_C(x).
\]
$s=0$ のときも $p_C(0)=0$ なので成立する．
劣加法性について，$a>p_C(x)$，$b>p_C(y)$ を取る．
$0\in C$ と凸性により，
$x\in aC$，$y\in bC$ が成り立つ．
実際，下限の定義から $x\in a'C$ となる $a'<a$ を取り，
$(a'/a)C\subset C$ を用いる．
従って
\[
 \frac{x+y}{a+b}
 =\frac a{a+b}\frac xa+\frac b{a+b}\frac yb\in C,
\]
ゆえに $p_C(x+y)\le a+b$ となる．
$a\downarrow p_C(x)$，$b\downarrow p_C(y)$ として
劣加法性を得る．

$x\in C$ なら $C$ が開であるから，
ある $\varepsilon>0$ に対して
$(1+\varepsilon)x\in C$ となる．
従って $p_C(x)\le1/(1+\varepsilon)<1$ である．
逆に $p_C(x)<1$ なら，
ある $t<1$ で $x\in tC$ となり，
$0\in C$ と凸性から $tC\subset C$ なので $x\in C$ である．
\end{proof}

\section{凸集合の分離}

\begin{theorem}[開凸集合とその外点の分離]\label{thm:open-sep}
$X$ を実ノルム空間，$C\subset X$ を空でない開凸集合，
$x_0\in X\setminus C$ とする．
このとき $0$ でない $f\in X^*$ が存在して
\[
 f(c)<f(x_0)\qquad(c\in C)
\]
が成り立つ．
\end{theorem}
\begin{proof}
$c_0\in C$ を固定し，
$U=C-c_0$ と $v=x_0-c_0$ を置く．
$U$ は $0$ を含む開凸集合，$v\notin U$ である．
$U$ のMinkowski汎関数を $p=p_U$ と書く．
補題\ref{lem:gauge}から $p(v)\ge1$ である．
$M=\mathbb Rv$ 上で
$g(tv)=t\,p(v)$ と定める．
$v\ne0$ なのでwell-definedで線形である．
$t\ge0$ なら
$g(tv)=p(tv)$，$t<0$ なら
$g(tv)=t\,p(v)\le0\le p(tv)$ である．
従って $g\le p|_M$ である．

実Hahn--Banachの定理により $g$ を
$F:X\to\mathbb R$ に拡張して $F\le p$ とできる．
$U$ が $B_X(0,r)$ を含むとすると
\[
 F(x)\le p(x)\le\frac{\|x\|}{r},
 \qquad
 -F(x)=F(-x)\le p(-x)\le\frac{\|x\|}{r}.
\]
よって $F$ は有界な実線形汎関数である．
$c\in C$ なら $c-c_0\in U$ なので
\[
 F(c-c_0)\le p(c-c_0)<1
 \le p(v)=F(v).
\]
従って $F(c)<F(x_0)$ となる．
また $F(v)=p(v)\ge1$ なので $F\ne0$ である．
$f=F$ として結論を得る．
\end{proof}

\begin{corollary}[開凸集合の相互分離]
実ノルム空間 $X$ において，
$A,B$ を互いに交わらない空でない凸集合とし，
$A$ は開であるとする．
このとき $0$ でない $f\in X^*$ が存在して
\[
 f(a)<f(b)\qquad(a\in A,\ b\in B)
\]
が成り立つ．
\end{corollary}
\begin{proof}
$U=A-B=\{a-b:a\in A,\ b\in B\}$ と置く．
$U$ は空でない凸集合であり，
$A$ が開であるから $U$ は開である．
また $A\cap B=\varnothing$ なので $0\notin U$ である．
定理\ref{thm:open-sep}を $U$ と $x_0=0$ に適用すると
ある $0\ne f\in X^*$ が存在して
$f(u)<f(0)=0$（$u\in U$）となる．
$u=a-b$ を代入すれば $f(a)<f(b)$ を得る．
\end{proof}

\begin{theorem}[閉凸集合と外点の強分離]\label{thm:strong-sep}
$X$ を実ノルム空間，
$C\subset X$ を空でない閉凸集合とし，
$x_0\notin C$ とする．
このとき $0\ne f\in X^*$ と実数 $\alpha$ が存在して
\[
 f(c)\le\alpha<f(x_0)\qquad(c\in C)
\]
が成り立つ．
\end{theorem}
\begin{proof}
$C$ は閉で $x_0\notin C$ だから
$d=d(x_0,C)>0$ である．
$0<\varepsilon<d$ を取り，
\[
 U=C+B_X(0,\varepsilon)
\]
と置く．$U$ は空でない開凸集合であり，
$\|x_0-c\|\ge d>\varepsilon$（$c\in C$）なので
$x_0\notin U$ である．
定理\ref{thm:open-sep}によって
$0\ne f\in X^*$ が存在して
\[
 f(c+b)<f(x_0)\qquad
 (c\in C,\ \|b\|<\varepsilon)
\]
となる．
固定した $c$ に対して $b$ に関する上限を取ると
\[
 f(c)+\varepsilon\|f\|\le f(x_0)
\]
を得る．実際，
$\sup_{\|b\|<\varepsilon}f(b)=\varepsilon\|f\|$ は，
実線形汎関数のノルムの定義と $b$ の符号の変更から従う．
従って
$\alpha=f(x_0)-\varepsilon\|f\|$
と置けば
$f(c)\le\alpha<f(x_0)$ が成り立つ．
\end{proof}

\begin{corollary}[閉凸集合の半空間表示]\label{cor:halfspaces}
実ノルム空間 $X$ の空でない閉凸集合 $C$ は，
$C$ を含む閉半空間の共通部分と一致する．
具体的には
\[
 C=\bigcap_{\substack{f\in X^*\\
       \sup_{c\in C}f(c)<\infty}}
 \left\{x\in X:
 f(x)\le \sup_{c\in C}f(c)\right\}
\]
が成り立つ．
\end{corollary}
\begin{proof}
右辺の各閉半空間は $C$ を含むから，
$C$ は右辺に含まれる．
逆に $x\notin C$ なら定理\ref{thm:strong-sep}から
ある $f\in X^*$ で
$\sup_{c\in C}f(c)<f(x)$ を満たすものが存在する．
特に $\sup_Cf$ は有限であり，
$x$ は対応する半空間に含まれない．
従って右辺は $C$ に含まれる．
\end{proof}

\begin{remark}
ここでの分離定理は実線形汎関数によって定式化した．
複素ノルム空間上では連続複素線形汎関数の実部を用いて
同様の分離を記述する．
すなわち不等式に現れる $f$ を
$\operatorname{Re}F$ と置き換える．
\end{remark}

\section{具体例とRieszの表現定理との比較}

\begin{example}[$\ell^1$ におけるノルムを与える汎関数]
$x=(x_n)\in\ell^1(\mathbb K)$，$x\ne0$ とする．
\[
 u_n=
 \begin{cases}
 \overline{x_n}/|x_n|& (x_n\ne0),\\
 0&(x_n=0)
 \end{cases}
\]
と置けば
$f_u(z)=\sum_{n=1}^\infty z_nu_n$ は
$\ell^1$ 上の連続線形汎関数であり，
$\|f_u\|=1$ および $f_u(x)=\|x\|_1$ を満たす．
\end{example}
\begin{proof}
$|u_n|\le1$ なので
$|f_u(z)|\le\sum_n|z_n|\le\|z\|_1$ となり，
$\|f_u\|\le1$ である．
定義から
$f_u(x)=\sum_n|x_n|=\|x\|_1$ なので，
$f_u$ のノルムは $1$ 以上である．
よって $\|f_u\|=1$ である．
\end{proof}

\begin{example}[評価汎関数]
$X=C([0,1],\mathbb K)$ に一様ノルムを入れる．
任意の $t_0\in[0,1]$ について
$\delta_{t_0}(f)=f(t_0)$ は
$\|\delta_{t_0}\|=1$ を満たす連続線形汎関数である．
\end{example}
\begin{proof}
$|\delta_{t_0}(f)|\le\|f\|_\infty$ なので
$\|\delta_{t_0}\|\le1$ である．
定数関数 $f\equiv1$ に対して
$\|f\|_\infty=1$，$|\delta_{t_0}(f)|=1$ だから
逆向きの不等式も成り立つ．
\end{proof}

\begin{remark}[Hahn--BanachとRieszの相違]
Hahn--Banachの定理は一般のノルム空間で有界線形汎関数の
\emph{拡張}を保証する．
他方，Hilbert空間のRieszの表現定理は
全ての有界線形汎関数が一つの内積
$x\mapsto\langle x,y\rangle$ により\emph{表現}されることを保証する．
後者はHilbert空間の内積構造を用いる別種の定理である．
\end{remark}

\section{結語}
Hahn--Banachの定理は，劣線形汎関数による支配条件を
保った線形汎関数の全空間への拡張を保証する．
この結果は，点を分離する連続線形汎関数の存在，
ノルムと距離の双対表示，
および凸集合の分離定理へと展開される．
これらは凸解析，最適化理論，弱位相の基礎をなす．

\section{到達課題の完全解答}
\noindent\textbf{Hahn--Banachの具体的な延長．}
\begin{proposition}
$X=\mathbb R^2$ にEuclidノルムを入れ，
$M=\{(x,0):x\in\mathbb R\}$ とし，
$f(x,0)=x$ と置く．
$f$ はノルム1であり，
$F(x,y)=x$ は $X$ へのノルム保存延長である．
\end{proposition}
\begin{proof}
$|f(x,0)|=|x|=\|(x,0)\|$ だから $\|f\|=1$．
$F|_M=f$ は明らかである．
任意の $(x,y)$ に対して
$|F(x,y)|=|x|\le\sqrt{x^2+y^2}$ だから
$\|F\|\le1$．$M$ 上でノルム1を達成するため
$\|F\|=1$ である．
\end{proof}
\begin{remark}
この例はHahn--Banachの存在定理を
具体的な公式で実現するが，一般のノルム空間で
延長が一意であることを主張しない．
\end{remark}
\noindent\textbf{到達確認．}
$M=\{(0,y)\}$，$f(0,y)=3y$ の
ノルム保存延長を一つ作れ．
\par\noindent\textbf{解答．}
$F(x,y)=3y$ であり，
$\|F\|=\|f\|=3$ となる．
\end{document}
